\subsection{固定分辨率碰撞矩轴的有理谱、Hankel 有限决定性与顶端剥离}\label{subsec:conclusion-fixedresolution-collision-hankel-prony-spectrum}
沿用定理 \ref{thm:pom-fiber-spectrum-resolvent-rational}、命题 \ref{prop:pom-power-sum-hankel-psd}、推论 \ref{cor:pom-crossq-logconvex-chain}、命题 \ref{prop:pom-escort-derivative-identities}、定理 \ref{thm:pom-fiber-spectrum-prony-hankel-rank}、定理 \ref{thm:pom-fiber-spectrum-prony-hankel-2r-reconstruction}、定理 \ref{thm:pom-fiber-spectrum-finite-reconstruction-sharp} 与定理 \ref{thm:xi-capacity-tail-histogram-moment-complete-statistic} 的记号。前文已经分别建立了固定分辨率碰撞矩的有理 resolvent、矩阶 Hankel 正半定、跨 \(q\) 对数凸性、escort 导数公式以及 \(2r\)-矩 Prony--Hankel 完全重构。将这些接口在同一固定分辨率层上闭合后，可把碰撞矩轴严格读作一个有限支撑离散谱的完整矩列；其有理生成函数、Hankel 秩、有限决定性、顶端谱恢复与连续容量曲线因而不再是平行描述，而成为同一有限谱几何的不同投影。

固定分辨率 \(m\ge 1\)。记纤维多重度的不同取值为
\[
1\le \delta_1<\delta_2<\cdots<\delta_{s_m},
\qquad
\mu_j:=\#\{x\in X_m:\ d_m(x)=\delta_j\}.
\]
于是对每个整数 \(q\ge 0\) 都有
\[
S_q(m)=\sum_{x\in X_m}d_m(x)^q=\sum_{j=1}^{s_m}\mu_j\delta_j^q.
\]

\begin{theorem}[固定分辨率碰撞矩轴的有理化与 Hankel 秩精确公式]\label{thm:conclusion-fixedresolution-collision-rational-hankel-rank}
\leanverified{paper\_conclusion\_fixedresolution\_collision\_rational\_hankel\_rank}
定义生成函数
\[
\mathcal F_m(z):=\sum_{q\ge 0}S_q(m)z^q
\]
与 Hankel 矩阵
\[
H_m^{(r)}:=\bigl(S_{i+j}(m)\bigr)_{0\le i,j\le r}.
\]
则成立：
\begin{enumerate}
  \item 在 \(|z|<\delta_{s_m}^{-1}\) 上有严格部分分式展开
  \[
  \mathcal F_m(z)=\sum_{j=1}^{s_m}\frac{\mu_j}{1-\delta_j z}.
  \]
  特别地，\(\mathcal F_m\) 为真有理函数，其极点恰为
  \[
  z=\delta_j^{-1}\qquad(1\le j\le s_m),
  \]
  并且全部为单极点。

  \item 对任意 \(r\ge 0\)，有精确秩公式
  \[
  \operatorname{rank}H_m^{(r)}=\min\{r+1,\ s_m\}.
  \]
  因而
  \[
  \det H_m^{(s_m-1)}>0,
  \qquad
  \det H_m^{(r)}=0\qquad(r\ge s_m).
  \]
\end{enumerate}
\end{theorem}
\begin{proof}
第 \(1\) 条只是定理 \ref{thm:pom-fiber-spectrum-resolvent-rational} 的记号重写：
\[
\mathcal F_m(z)
=
\sum_{q\ge 0}\sum_{j=1}^{s_m}\mu_j(\delta_j z)^q
=
\sum_{j=1}^{s_m}\mu_j\sum_{q\ge 0}(\delta_j z)^q
=
\sum_{j=1}^{s_m}\frac{\mu_j}{1-\delta_j z},
\]
其中交换求和在 \(|z|<\delta_{s_m}^{-1}\) 上合法。由于 \(\delta_j\) 两两不同，故极点互异且皆为一阶。

第 \(2\) 条可直接写成 Vandermonde 分解。令
\[
V_r:=\bigl(\delta_j^i\bigr)_{1\le j\le s_m,\ 0\le i\le r},
\qquad
D:=\mathrm{diag}(\mu_1,\dots,\mu_{s_m}).
\]
则
\[
(H_m^{(r)})_{ij}=S_{i+j}(m)=\sum_{k=1}^{s_m}\mu_k\delta_k^{i+j},
\]
故
\[
H_m^{(r)}=V_r^{\!\top}DV_r.
\]
由于 \(D\) 正定，故 \(\operatorname{rank}H_m^{(r)}=\operatorname{rank}V_r\)。而 \(V_r\) 为 Vandermonde 型矩阵，节点 \(\delta_j\) 两两不同，所以其秩恰为 \(\min\{r+1,s_m\}\)。遂得所述秩公式与行列式结论。证毕。
\end{proof}

\begin{corollary}[首 \texorpdfstring{$2s_m$}{2sm} 阶矩的全谱充分性]\label{cor:conclusion-fixedresolution-first-2sm-sufficiency}
有限数据
\[
S_0(m),S_1(m),\dots,S_{2s_m-1}(m)
\]
唯一决定多重集
\[
\{(\delta_j,\mu_j)\}_{j=1}^{s_m},
\]
从而唯一决定：
\[
\{S_q(m)\}_{q\ge 0},
\qquad
N_m(t):=\#\{x\in X_m:\ d_m(x)\ge t\},
\qquad
\mathcal C_m^{\mathrm{cont}}(\cdot).
\]
因此，在固定分辨率上，连续容量曲线与完整矩谱之间的充分统计等价可 sharpen 为：首 \(2s_m\) 阶碰撞矩已经是整条有限谱的全统计量。
\end{corollary}
\begin{proof}
由定理 \ref{thm:conclusion-fixedresolution-collision-rational-hankel-rank} 的秩公式，固定 \(m\) 时最小 Hankel 秩恰为 \(s_m\)。于是定理 \ref{thm:pom-fiber-spectrum-prony-hankel-2r-reconstruction} 与定理 \ref{thm:pom-fiber-spectrum-finite-reconstruction-sharp} 直接给出：首 \(2s_m\) 阶矩唯一恢复全部原子 \(\delta_j\) 及其权重 \(\mu_j\)，并且该阈值在一般意义下不可降低。

一旦 \(\{(\delta_j,\mu_j)\}\) 已知，则
\[
S_q(m)=\sum_{j=1}^{s_m}\mu_j\delta_j^q
\]
对任意 \(q\ge 0\) 立即确定；而
\[
N_m(t)=\sum_{\delta_j\ge t}\mu_j
\]
亦随之唯一恢复。最后由定理 \ref{thm:xi-capacity-tail-histogram-moment-complete-statistic}，尾计数函数 \(N_m\)、连续容量曲线 \(\mathcal C_m^{\mathrm{cont}}\) 与全矩谱彼此唯一决定，故结论成立。证毕。
\end{proof}

\begin{theorem}[连续矩比流的顶端谱恢复与次顶端缺口律]\label{thm:conclusion-fixedresolution-moment-ratio-topgap}
定义相邻矩比
\[
R_q(m):=\frac{S_{q+1}(m)}{S_q(m)}
\qquad(q\ge 0).
\]
则成立：
\begin{enumerate}
  \item 序列 \(q\mapsto R_q(m)\) 单调不减。

  \item 若 \(s_m\ge 2\)，则对每个 \(q\ge 0\) 有
  \[
  \delta_1<R_q(m)<\delta_{s_m},
  \]
  并且
  \[
  \lim_{q\to\infty}R_q(m)=\delta_{s_m}.
  \]

  \item 若 \(s_m\ge 2\)，则存在精确指数缺口律
  \[
  \lim_{q\to\infty}\frac1q\log\bigl(\delta_{s_m}-R_q(m)\bigr)
  =
  \log\frac{\delta_{s_m-1}}{\delta_{s_m}}.
  \]
\end{enumerate}
特别地，仅由相邻矩比序列 \(\{R_q(m)\}_{q\ge 0}\) 即可恢复顶端双谱：
\[
\delta_{s_m}=\lim_{q\to\infty}R_q(m),
\]
\[
\delta_{s_m-1}
=
\delta_{s_m}\cdot
\exp\!\left(
\lim_{q\to\infty}\frac1q\log(\delta_{s_m}-R_q(m))
\right),
\]
并且
\[
\mu_{s_m}
=
\lim_{q\to\infty}\frac{S_q(m)}{\delta_{s_m}^q}.
\]
\end{theorem}
\begin{proof}
由推论 \ref{cor:pom-crossq-logconvex-chain}，
\[
S_{q+1}(m)^2\le S_q(m)S_{q+2}(m),
\]
故 \(R_q(m)\le R_{q+1}(m)\)，即矩比流单调不减。

另一方面，
\[
R_q(m)
=
\frac{\sum_{j=1}^{s_m}\mu_j\delta_j^{q+1}}{\sum_{j=1}^{s_m}\mu_j\delta_j^q}
=
\sum_{j=1}^{s_m}\pi_{j,q}\delta_j,
\qquad
\pi_{j,q}:=\frac{\mu_j\delta_j^q}{S_q(m)}.
\]
因此 \(R_q(m)\) 是 \(\{\delta_j\}\) 的凸组合。若 \(s_m\ge 2\)，则每个 \(\pi_{j,q}\) 都严格为正，故
\[
\delta_1<R_q(m)<\delta_{s_m}.
\]
又因为
\[
\pi_{s_m,q}
=
\frac{\mu_{s_m}\delta_{s_m}^q}{\sum_{j=1}^{s_m}\mu_j\delta_j^q}
\longrightarrow 1,
\]
从而 \(R_q(m)\to \delta_{s_m}\)。

再写顶端缺口：
\[
\delta_{s_m}-R_q(m)
=
\frac{\sum_{j=1}^{s_m-1}\mu_j\delta_j^q(\delta_{s_m}-\delta_j)}
{\mu_{s_m}\delta_{s_m}^q+\sum_{j=1}^{s_m-1}\mu_j\delta_j^q}.
\]
分子与分母中均由次顶层 \(j=s_m-1\) 主导，故
\[
\delta_{s_m}-R_q(m)
=
\frac{\mu_{s_m-1}(\delta_{s_m}-\delta_{s_m-1})}{\mu_{s_m}}
\left(\frac{\delta_{s_m-1}}{\delta_{s_m}}\right)^q
\bigl(1+o(1)\bigr),
\]
从而
\[
\lim_{q\to\infty}\frac1q\log\bigl(\delta_{s_m}-R_q(m)\bigr)
=
\log\frac{\delta_{s_m-1}}{\delta_{s_m}}.
\]
前两条恢复公式由此直接得到。最后，
\[
\frac{S_q(m)}{\delta_{s_m}^q}
=
\mu_{s_m}+\sum_{j=1}^{s_m-1}\mu_j\left(\frac{\delta_j}{\delta_{s_m}}\right)^q
\longrightarrow \mu_{s_m},
\]
故顶层权重亦可由尾部恢复。证毕。
\end{proof}

\begin{theorem}[有限分辨率碰撞谱的逐层顶端剥离]\label{thm:conclusion-fixedresolution-topdown-peeling}
\leanverified{paper\_conclusion\_fixedresolution\_topdown\_peeling}
定义递归残差矩列
\[
T_q^{(0)}(m):=S_q(m).
\]
对 \(r=0,1,\dots,s_m-1\)，递归定义
\[
\delta_{s_m-r}
:=
\lim_{q\to\infty}\bigl(T_q^{(r)}(m)\bigr)^{1/q},
\qquad
\mu_{s_m-r}
:=
\lim_{q\to\infty}\frac{T_q^{(r)}(m)}{\delta_{s_m-r}^q},
\]
\[
T_q^{(r+1)}(m):=
T_q^{(r)}(m)-\mu_{s_m-r}\delta_{s_m-r}^q.
\]
则上述极限全部存在，并且恰按顶端顺序逐层恢复
\[
(\delta_{s_m},\mu_{s_m}),\ (\delta_{s_m-1},\mu_{s_m-1}),\ \dots,\ (\delta_1,\mu_1).
\]
终端时有
\[
T_q^{(s_m)}(m)\equiv 0.
\]
换言之，固定分辨率的整个碰撞谱都可由正矩尾部作 top-down peeling 而被完全层析。
\end{theorem}
\begin{proof}
当 \(r=0\) 时，
\[
T_q^{(0)}(m)=\sum_{j=1}^{s_m}\mu_j\delta_j^q,
\]
故
\[
\lim_{q\to\infty}\bigl(T_q^{(0)}(m)\bigr)^{1/q}=\delta_{s_m},
\qquad
\lim_{q\to\infty}\frac{T_q^{(0)}(m)}{\delta_{s_m}^q}=\mu_{s_m}.
\]
假设归纳地已经得到
\[
T_q^{(r)}(m)=\sum_{j=1}^{s_m-r}\mu_j\delta_j^q.
\]
则与 \(r=0\) 完全相同的最大项主导给出
\[
\lim_{q\to\infty}\bigl(T_q^{(r)}(m)\bigr)^{1/q}=\delta_{s_m-r},
\qquad
\lim_{q\to\infty}\frac{T_q^{(r)}(m)}{\delta_{s_m-r}^q}=\mu_{s_m-r}.
\]
减去该顶层项后得到
\[
T_q^{(r+1)}(m)=\sum_{j=1}^{s_m-r-1}\mu_j\delta_j^q,
\]
故归纳闭合。最终所有谱点全部剥离，因而 \(T_q^{(s_m)}(m)\equiv 0\)。证毕。
\end{proof}

\begin{proposition}[碰撞矩轴是一条有限指数族的对数配分函数]\label{prop:conclusion-fixedresolution-logpartition-escort-convexity}
对实参数 \(q\in\RR\) 定义
\[
F_m(q):=\log\Bigl(\sum_{x\in X_m}d_m(x)^q\Bigr)=\log S_q(m),
\qquad
\nu_{m,q}(x):=\frac{d_m(x)^q}{S_q(m)}.
\]
则
\[
F_m'(q)=\sum_{x\in X_m}\nu_{m,q}(x)\log d_m(x),
\]
\[
F_m''(q)=\operatorname{Var}_{\nu_{m,q}}\!\bigl(\log d_m(x)\bigr)\ge 0.
\]
因此 \(F_m\) 为凸函数，并且
\[
F_m\ \text{严格凸}
\iff
\#\{d_m(x):x\in X_m\}\ge 2.
\]
\end{proposition}
\begin{proof}
这与命题 \ref{prop:pom-escort-derivative-identities} 同型，但此处可直接在固定 \(m\) 下计算：因为
\[
S_q(m)=\sum_{x\in X_m}e^{q\log d_m(x)}
\]
是有限个指数函数之和，故对全部实数 \(q\) 都可微，并且
\[
F_m'(q)
=
\frac{\sum_x d_m(x)^q\log d_m(x)}{S_q(m)}
=
\sum_x \nu_{m,q}(x)\log d_m(x).
\]
再求一次导数即得
\[
F_m''(q)
=
\frac{\sum_x d_m(x)^q(\log d_m(x))^2}{S_q(m)}
-
\left(
\frac{\sum_x d_m(x)^q\log d_m(x)}{S_q(m)}
\right)^2,
\]
即为题述方差公式。

若所有 \(d_m(x)\) 相等，则 \(\log d_m(x)\) 在 \(\nu_{m,q}\) 下为常数，故 \(F_m''\equiv 0\)。反之，只要至少有两个不同多重度值，\(\nu_{m,q}\) 对它们都赋正质量，于是该方差对每个 \(q\in\RR\) 都严格为正，故 \(F_m\) 严格凸。证毕。
\end{proof}
\leanverified{paper\_conclusion\_fixedresolution\_logpartition\_escort\_convexity}

\begin{theorem}[主理想型 Hankel-null 的精确有限决定性]\label{thm:conclusion-fixedresolution-principal-hankelnull-finite-determinacy}
设固定分辨率 \(m\) 的碰撞矩轴具有最小消去多项式
\[
P_m(\lambda)=\lambda^d+a_1\lambda^{d-1}+\cdots+a_d\in\ZZ[\lambda],
\]
并且对应 Hankel-null 理想由 \(P_m\) 主生成。则：
\begin{enumerate}
  \item 矩轴满足整系数递推
  \[
  S_{q+d}(m)+a_1S_{q+d-1}(m)+\cdots+a_dS_q(m)=0
  \qquad(q\ge 0).
  \]

  \item 数据
  \[
  P_m(\lambda)
  \qquad\text{与}\qquad
  S_0(m),\dots,S_{d-1}(m)
  \]
  唯一决定整个序列 \(\{S_q(m)\}_{q\ge 0}\)。

  \item 若仅知最小 Hankel 秩等于 \(d\)，则首 \(2d\) 阶矩
  \[
  S_0(m),\dots,S_{2d-1}(m)
  \]
  已足以同时恢复 \(P_m\) 与全部高阶矩。
\end{enumerate}
\end{theorem}
\begin{proof}
由定理 \ref{thm:pom-fiber-spectrum-prony-hankel-rank}，最小消去多项式的次数等于最小 Hankel 秩；因此在当前固定分辨率情形中 \(d=s_m\)。又因为不同谱点恰为 \(\delta_1,\dots,\delta_d\)，所以
\[
P_m(\lambda)=\prod_{j=1}^{d}(\lambda-\delta_j).
\]
对每个 \(j\) 与每个 \(q\ge 0\) 都有
\[
\delta_j^{q+d}+a_1\delta_j^{q+d-1}+\cdots+a_d\delta_j^q=0.
\]
乘以 \(\mu_j\) 并对 \(j\) 求和，即得第 \(1\) 条。

第 \(2\) 条是常系数线性递推的唯一延拓性：一旦递推式与前 \(d\) 项已知，所有后续项都被逐步唯一决定。

第 \(3\) 条则正是定理 \ref{thm:pom-fiber-spectrum-prony-hankel-2r-reconstruction} 在 \(r=d\) 时的直接应用；同时定理 \ref{thm:pom-fiber-spectrum-finite-reconstruction-sharp} 说明这一 \(2d\) 阈值在一般意义下是尖锐的。证毕。
\end{proof}

\begin{theorem}[统一 Hankel 秩界等价于有限谱字母表]\label{thm:conclusion-fixedresolution-uniform-hankel-rank-alphabet}
设 \(\mathcal M\subset\NN\) 为一族分辨率。对每个 \(m\in\mathcal M\)，记
\[
s_m:=\#\{d_m(x):x\in X_m\}.
\]
则对任意固定整数 \(D\ge 1\)，下列命题等价：
\begin{enumerate}
  \item 对所有 \(m\in\mathcal M\) 都有
  \[
  s_m\le D.
  \]

  \item 对所有 \(m\in\mathcal M\) 与所有 \(r\ge 0\)，有
  \[
  \operatorname{rank}H_m^{(r)}\le D.
  \]

  \item 对所有 \(m\in\mathcal M\)，都有
  \[
  \det H_m^{(D)}=0,
  \]
  等价地，每个固定分辨率上的 \((D+1)\times(D+1)\) Hankel 主子式都消失。

  \item 对所有 \(m\in\mathcal M\)，碰撞矩轴都满足阶数至多 \(D\) 的线性递推。
\end{enumerate}
换言之，跨分辨率存在统一 Hankel 秩界，当且仅当整个分辨率族在 \(q\)-轴上只使用一个固定大小的有限谱字母表。
\end{theorem}
\begin{proof}
由定理 \ref{thm:conclusion-fixedresolution-collision-rational-hankel-rank}，
\[
\operatorname{rank}H_m^{(r)}=\min\{r+1,s_m\},
\]
故第 \(1\) 条与第 \(2\) 条立即等价。将 \(r=D\) 代入可知第 \(1\) 条等价于
\[
\operatorname{rank}H_m^{(D)}\le D,
\]
这又等价于 \(\det H_m^{(D)}=0\)，故第 \(1\)、第 \(2\)、第 \(3\) 条等价。

另一方面，由定理 \ref{thm:conclusion-fixedresolution-principal-hankelnull-finite-determinacy}，固定 \(m\) 时最小递推阶数恰等于最小 Hankel 秩，也恰等于 \(s_m\)。因此“所有 \(m\) 的递推阶数都至多为 \(D\)”与“所有 \(m\) 的 \(s_m\) 都至多为 \(D\)”完全等价，从而第 \(4\) 条亦与前述三条等价。证毕。
\end{proof}

\begin{conjecture}[genuine Fold 的跨分辨率 Hankel 秩无界]\label{conj:conclusion-fixedresolution-fold-unbounded-hankel-rank}
对真实 Fold 族 \((\Fold_m)_{m\ge 1}\)，应有
\[
\sup_{m\ge 1}s_m=+\infty.
\]
等价地，对任意固定 \(D\ge 1\)，总存在充分大的 \(m\) 使得
\[
\det H_m^{(D)}\neq 0.
\]
若该猜想成立，则碰撞矩轴的 finite determinacy 只能是逐分辨率的，而不可能由任何固定阶 Hankel 证书在全部分辨率上一劳永逸地封闭；moment-kernel 层级的持续增加也就不是技术冗余，而是由真实 Fold 的谱字母表持续增生所强制。
\end{conjecture}

这条固定分辨率链可压缩为
\[
\boxed{
\text{有限支撑碰撞谱}
\Longleftrightarrow
\text{有理生成函数}
\Longleftrightarrow
\text{Hankel 秩 } s_m
\Longleftrightarrow
\text{首 }2s_m\text{ 阶矩的完全决定性}
}
\]
\[
\boxed{
\text{相邻矩比流}
\Longrightarrow
\text{顶端双谱与权重恢复}
\Longrightarrow
\text{top-down peeling}
\Longrightarrow
\text{连续容量曲线的有限谱层析}
}
\]
因此，固定分辨率的碰撞矩轴并不是一条抽象的高阶 Rényi 曲线，而是一条可被 Hankel--Prony 机制完全层析的有限离散谱；其真正承载的不是单一标量熵率，而是整条可审计、可剥离、可反演的有限谱几何。
